\section{Conclusion}

This project built a working prototype for decentralized LLM service discovery
and routing. Each node can host or access a language model, advertise
capabilities, discover peers through libp2p and a Kademlia DHT, classify
queries, rank candidates, and route execution. The system also includes HTTP
APIs, DIAL Chat, and a benchmark suite for latency, scale, dropout, and policy
weights.

The main result is that a peer-to-peer layer is a credible starting point
for decentralized AI services when its role is clear. In local benchmarks,
clean and moderate-latency networks route reliably, mixed-capability queries
work well in clean conditions, and the capability vector model gives a path
toward richer matching. The evaluation also clarifies that the aim is not to
route every request to the same strongest node, but to choose suitable peers
while keeping the network robust. The hard part is clear: candidate visibility
does not yet scale well enough, and dropout, stale availability information,
and timeout-heavy failure detection remain important bottlenecks. This points
first to DHT/provider-discovery tuning, faster failed-peer handling, and
bounded recommendation expansion. Load-aware routing remains a natural next
direction once discovery is stronger.

The broader vision is an open AI service network where providers can join,
advertise specialized models, build reputation, and be selected by transparent
matching policies. Work on LLM routing, agent discovery, PPAI-style delegation,
verifier models, and privacy-preserving routing suggests that this is becoming
an active research direction. Blockchain is not needed for first-layer
discovery, but may help later with payment settlement and verifiable
reputation. In that sense, DIAL is best understood as a lightweight P2P
foundation for discovery and routing today, with a path toward incentives,
trust, and multi-model collaboration later.
